\documentclass[11pt,a4paper]{article}
\usepackage[T1]{fontenc}
\usepackage[utf8]{inputenc}
\usepackage{lmodern}
\usepackage[a4paper,hmargin=22mm,vmargin=20mm]{geometry}
\usepackage{amsmath,amssymb,amsthm,booktabs,tabularx,array}
\usepackage{xurl,xcolor}
\usepackage[colorlinks,allcolors=blue!55!black]{hyperref}
\renewcommand{\abstractname}{Resumen}
\renewcommand{\proofname}{Demostración}
\renewcommand{\refname}{Referencias}
\newtheorem{theorem}{Teorema}
\newcommand{\bits}{\{0,1\}}
\newcommand{\hh}{H_{\min}}
\newcommand{\AES}{\operatorname{AES256}}
\newcommand{\SHA}{\operatorname{SHA256}}
\newcommand{\TV}{\operatorname{TV}}
\newcommand{\inc}{\operatorname{inc}}
\newcolumntype{Y}{>{\raggedright\arraybackslash}X}
\title{DRG.3 y conservación ITS\\
\large Tupla AIS v3.0, evidencia algorítmica y un perfil exacto de CTR\_DRBG}
\author{Addendum técnico para Jordi Prieto Gallego}
\date{20 de noviembre de 2026}

\begin{document}
\maketitle
\begin{abstract}
Se da la tupla completa del ejemplo AES/SHA de AIS 20/31 v3.0 y se separa
su prueba exacta de conservación ITS de su prueba de seguridad
computacional en un modelo ideal. Un argumento por requisito identifica
la evidencia que puede aportarse para PTG.3.2. Se demuestra además que
CTR\_DRBG con AES-256 sin función de derivación puede preservar
exactamente los 128 bits de entrada fresca por bloque si se inyectan
en el campo $V$ mediante el \emph{additional input} de 384 bits.
El argumento usa el algoritmo de generación de SP 800-90A Rev. 1
y tiene un alcance separado del análisis DRG.3 de CTR\_DRBG.
\end{abstract}

\section{Tupla exacta y mapa ITS del ejemplo de AIS}

Sea $p\in\bits^{128}$ una constante pública. Sea $\iota(a)$ el
padding de $a$ con ceros a la derecha hasta 128 bits. Las componentes
de los §§996--997, completando $I$ con el límite de §976, son
\[
\begin{aligned}
\mathcal T&=(\mathcal S,\mathcal S_{\rm req},\mathcal R,\mathcal A,
 I,\phi,\phi_{\rm req},\phi_0,\psi),\\
\mathcal S&=\bits^{256}, &
\mathcal S_{\rm req}&=\bits^{256}\times\bits^{128},\\
\mathcal R&=\bits^{128}, &
\mathcal A&=\bigcup_{\ell=0}^{128}\bits^\ell,\\
I&=\{1,\ldots,128\}, &
\phi_{\rm req}(s,a)&=(s,\iota(a)),\\
\phi_0(s,u)&=(s,u), &
\psi(s,u)&=\AES_s(p\oplus u),\\
\phi(s,a)&=\SHA\big(s\mathbin\|(p\oplus\iota(a))\big).
\end{aligned}
\]
El núcleo produce un número interno por petición. Si se solicitan menos
de 128 bits se entrega la proyección especificada por el perfil.
El perfil PTG.3 de la nota principal toma $I=\{128\}$ y una sola
entrada fresca $a\in\bits^{128}$ por llamada. El caso autónomo
$a=\varnothing$ es parte del análisis algorítmico. En el perfil fijo
$\iota(a)=a$. Para entradas de longitud variable, $\iota$ puede identificar
cadenas distintas, como la cadena vacía y una cadena de ceros
\cite[§§127--147, 976 y 996--998]{ais}.

\begin{theorem}[Conservación por bloque]
Sea $Z$ la información lateral, que puede incluir el estado actual $S$.
Para $Y=\AES_S(p\oplus A)$ y $A\in\bits^{128}$,
\[
\begin{aligned}
\hh(Y\mid S,Z)&=\hh(A\mid S,Z),\\
\TV(P_{Y,S,Z},U_{128}\otimes P_{S,Z})
  &=\TV(P_{A,S,Z},U_{128}\otimes P_{S,Z}).
\end{aligned}
\]
\end{theorem}
\begin{proof}
Para cada estado $s$ la inversa es
$a=p\oplus\AES_s^{-1}(y)$. Esto permuta las probabilidades de
$A$ en cada contexto $(s,z)$. Preserva el máximo de las
probabilidades y la distancia a la distribución uniforme.
Con $Z$ cuántico y $S$ clásico, la permutación controlada por $S$
es unitaria, preserva la distancia de traza y permite relabelar
de forma reversible las respuestas de una estrategia de adivinación.
\end{proof}

La biyección actúa sobre $a\mapsto y$ para cada $s$. La transición
$s\mapsto\phi(s,a)$ para $a$ conocido tiene un requisito
computacional de unidireccionalidad. Ambas propiedades operan en
mapas distintos y son compatibles. La entropía acreditada para PTG.3
proviene del $A$ físico. Si $a$ se fija a vacío y sólo queda una
semilla inicial de 256 bits, el soporte de $q$ salidas completas
tiene tamaño máximo $2^{256}$ y para $q>2$ la distancia a uniforme
es al menos $1-2^{256-128q}$. La nota principal da la prueba de
secuencia con entradas frescas y la cota para salidas truncadas
\cite[§§119, 390(iii), 397(ii), 402 y 996--998]{ais}.

\section{Demostración computacional en un modelo ideal}

\paragraph{Juego.}
Se estudia una época de $N$ peticiones. Para este argumento, $A_i$ puede
ser vacío, público o elegido adaptativamente por el adversario a partir de
lo observado. Se permite elegir $b_i=p\oplus\iota(A_i)$ arbitrario
de 128 bits. El atacante conoce el diseño y puede consultar el
cifrado ideal AES de 256 bits de clave y 128 bits de bloque, en
ambas direcciones, y un oráculo aleatorio de 256 bits para SHA
restringido a entradas de 384 bits. Los dos oráculos son
independientes. Cada petición aplica
$Y_i=\mathcal E_{S_i}(b_i)$ y
$S_{i+1}=\mathcal H(S_i\|b_i)$ para $i=0,\ldots,N-1$.
El objetivo es adivinar un bloque oculto a partir del transcript
de otros bloques, antes o después del objetivo según el juego.
Para enhanced backward secrecy también se entrega $S_N$ tras
completar la última transición y se dan por conocidos los $b_i$.
Puede haber información lateral clásica previa a la época.
Se asume que el máximo puntual de la probabilidad de cada valor
de $S_0$ condicionado a dicha información es $\le2^{-h}$ y que
la siembra no depende de los oráculos ideales.

\begin{theorem}[Ventaja en el modelo ideal]
Con a lo sumo $Q$ consultas externas a ambos oráculos en total,
y sin exponer estados intermedios del horizonte, la diferencia
respecto a un transcript con bloques uniformes independientes
está acotada, para los juegos indicados, por
\[
\delta\,=\,\min\left(1, \frac{N(N+1)}{2\cdot2^{256}}
 +Q\big[2^{-h}+(N-1)2^{-256}\big]\right).
\]
En particular, la probabilidad de acertar un bloque oculto es
a lo sumo $2^{-128}+\delta$.
\end{theorem}
\begin{proof}
En el mundo ideal de comparación se muestrean $S_1,\ldots,S_N$
como etiquetas uniformes independientes y se usan $N$ bloques
de salida uniformes e independientes para construir el
transcript. Los oráculos externos se responden con tablas
independientes. Se acopla con el generador real idealizado
programando internamente $\mathcal H(S_i\|b_i)=S_{i+1}$ y
$\mathcal E_{S_i}(b_i)=Y_i$. Mientras las etiquetas de estado sean
distintas y ninguna consulta externa use una clave oculta $S_i$
o una entrada SHA con ese prefijo de estado, cada consulta
interna SHA llega a un punto nuevo y cada evaluación AES es la
primera bajo su clave. Sus respuestas tienen la distribución
muestreada. También se mantiene el acoplamiento con $S_N$
expuesto porque ese estado terminal queda fuera del conjunto
de claves ocultas y nunca genera un bloque de la época analizada.

En el mundo de comparación, la probabilidad de colisión de
cualquiera de los $N+1$ estados es como máximo
$\binom{N+1}{2}2^{-256}$, incluso con $S_0$ no uniforme.
Como las salidas y las respuestas externas están simuladas
independientemente de las etiquetas ocultas, cada consulta
sigue teniendo probabilidad $\le2^{-h}$ de acertar $S_0$
y $\le2^{-256}$ de acertar cada otro estado oculto. La unión
de los eventos cubre la primera divergencia, también cuando
las elecciones de $b_i$ y las consultas son adaptativas.
En el transcript de comparación el bloque reservado es
uniforme dado lo observado. Esto da la cota de adivinación.
\end{proof}

\paragraph{Dimensión del estado en el modelo ideal.}
Póngase $M=2^{128}$ y elíjanse cinco bloques distintos $b_0,\ldots,b_4$.
Las firmas de un paso $(\mathcal E_s(b_j))_{j=0}^4$
son independientes entre claves y uniformes sobre $M(M-1)\cdots(M-4)$
resultados posibles. La probabilidad de que exista una pareja de
claves con la misma firma está acotada por
$\binom{2^{256}}{2}/[M(M-1)\cdots(M-4)]$, del orden de $2^{-129}$.
Así, con la probabilidad indicada, ninguna representación determinista
de menos de 256 bits puede reproducir la respuesta de un paso para todos
los estados y todas las entradas. Las firmas caracterizan el mapa de
respuestas de un estado matemático y no conceden clonación al atacante.
El juego anterior aporta la parte complementaria: observaciones reales
de una época no permiten resolver prácticamente el estado en el modelo.
La dimensión estructural y la entropía de siembra se acreditan por separado.

La prueba modela SHA y AES como dos objetos ideales y usa la
frontera de una época. Por ejemplo, que AES sea una PRP de clave
fija no implica por sí solo el comportamiento en una trayectoria
de claves derivadas. La hipótesis $h$ es una cota de
adivinación del estado inicial, no la dimensión del estado
eficaz. La exposición se sitúa entre peticiones. La
implementación debe modelar borrado, copias de claves,
concurrencia y exposición durante una petición. Con un bloque
de salida, la referencia de adivinación es $2^{-128}$.
Para fijar $N$ y los parámetros de evaluación se usa la
fuerza de seguridad, dimensión y vida útil requeridas por AIS
\cite[§§93--103, 106--125, 195--198 y 211--214]{ais}.

\section{Qué se demuestra por requisito y qué falta para AIS}

\noindent\small
\begin{tabularx}{\linewidth}{@{}lYY@{}}
\toprule
Requisito & Argumento algorítmico & Evidencia que debe aportar la instancia\\
\midrule
DRG.3.3 &
$S$ tiene 256 bits y toda la clave interviene en AES y en SHA.
En el juego ideal, la trayectoria desde semilla uniforme oculta
no ofrece un atajo para recuperar bits de la clave salvo
acertar etiquetas o sufrir colisiones. &
Justificar al menos 246 bits de estado eficaz en el modelo de
§§93--103 con salidas conocidas, apego a la tupla, parámetros
de seedlife y ausencia de representación reducida. \\
DRG.3.5 &
El juego con el bloque oculto después de los observados
da forward secrecy de números internos. &
Documentar exposición, longitud, límites de peticiones, forma
de agregación y supuestos computacionales para AES/SHA.\\
DRG.3.6 &
El juego con el bloque oculto antes de los observados
da backward secrecy de números internos. &
Incluir en el análisis las vistas permitidas de salidas futuras
y cómo se gestionan instancias y épocas.\\
DRG.3.7 &
El juego con $S_N$ comprometido y entradas conocidas cubre
enhanced backward secrecy para bloques previos de la época. &
Seleccionar granularidad de número interno. Comprobar
transición antes de entrega, borrado del estado anterior
y hacer corresponder el compromiso real con esa frontera.\\
DRG.3.8 &
El juego permite $b_i$ arbitrarios, repetidos y adaptativos.
La cota ideal no requiere frescura del additional input. &
Argumentar por qué la composición AES/SHA sostiene esa
resistencia realista, y comprobar que aliases de padding,
errores, resets y rollback no introducen otras transiciones.\\
DRG.3.9 &
Funciones criptográficas reconocidas. La transición ideal
es unidireccional para una entrada adicional conocida y
estado oculto de alta entropía. &
Citar §998, precisar el uso de la hipótesis de preimagen y
mostrar conformidad de las primitivas y codificación.\\
\bottomrule
\end{tabularx}
\normalsize

\paragraph{Suficiencia para un expediente.}
BSI afirma conformidad algorítmica del ejemplo en §998 y permite
invocar propiedades aceptadas de primitivas reconocidas en §§99--100.
Citar el ejemplo, acreditar equivalencia del núcleo y aportar la matriz
anterior constituye una base directa para PTG.3.2, cuyo §380 remite
expresamente a DRG.3.3 y DRG.3.5--8. La prueba ideal refuerza el
argumento y hace sus hipótesis revisables. La aceptación de que dichas
hipótesis y la implementación cubren el TOE corresponde al evaluador y
al organismo. La 9-tupla por sí sola especifica las funciones y el flujo.
La nota no da por aceptado un producto ni presenta el modelo ideal como
un teorema sobre AES y SHA concretos.

Para un claim DRG.3 autónomo completo, documentar también DRG.3.1
(origen de semilla), DRG.3.2 (límite $2^{19}$ bits por petición,
satisfecho con 128), DRG.3.4 (entropía del estado eficaz tras
siembra/reseed) y DRG.3.10 (inconspicuidad), además de DRG.3.9.
La 4-tupla de siembra propuesta es
$(\bits^{256},\{\varnothing\},\mathcal S,\phi_{\rm seed})$
con $\phi_{\rm seed}(M,\varnothing)=M$. La de reseed es
$(\bits^{256},\{\varnothing\},\mathcal S,\phi_{\rm reseed})$
con $\phi_{\rm reseed}(s,M,\varnothing)=M$. $M$ lo forman
dos bloques PTG.2 dedicados. La cota puntual de $128\cdot0.98$
para cada bloque condicionado al historial da $250.88$ bits
para $M$, si la fuente demuestra esa hipótesis
\cite[§§146--164, 207, 211--214, 380 y 996--998]{ais}.

\section{CTR\_DRBG: dónde se conserva ITS}

\subsection{El núcleo sin función de derivación}

Se considera AES-256, $\mathrm{seedlen}=384$, $\mathrm{blocklen}=128$,
y contador de 128 bits. Escribimos $s=(K,V)$ y omitimos de las ecuaciones
el contador de reseed, que gobierna cuándo se permite generar.
Sean $\inc(v)=v+1\bmod 2^{128}$ y
\[
G(s)=
 \AES_K(\inc(V))\,\|\,\AES_K(\inc^2(V))\,\|\,\AES_K(\inc^3(V))
 =G_K(s)\|G_V(s),
\]
donde los subíndices denotan las proyecciones de 256 y 128 bits,
respectivamente. El \emph{Update} de SP 800-90A Rev. 1 calcula
\[
\operatorname{Update}(D,s)=G(s)\oplus D,\qquad D\in\bits^{384}.
\]
La igualdad interpreta los 384 bits de resultado como la pareja
$(K',V')$. Fijado $s$, este mapa es biyectivo en $D$.
En \emph{Instantiate} sin derivación, $s$ es cero y $D$ es el
XOR de la entrada de entropía con la personalización rellenada.
En \emph{Reseed} sin derivación, $s$ es el estado antiguo y $D$
es el XOR de la entrada de entropía con el additional input
rellenado. Fijados los respectivos datos auxiliares, ambos pasos
preservan exactamente la entropía de la entrada de entropía
de 384 bits \emph{en el estado resultante}.
El estándar exige entrada \emph{full entropy} para esos
pasos sin derivación \cite[§10.2.1, tabla 3 y §§10.2.1.2--10.2.1.4]{nist}.

\subsection{Una forma de preservar ITS en la salida estándar}

Sea $A\in\bits^{128}$ la salida fresca del PTG.2 y tome la llamada
a \emph{Generate} de 128 bits con el siguiente argumento presente:
\[
 \boxed{D=0^{256}\|A\quad\text{como additional input de 384 bits.}}
\]
Incluso cuando $A=0^{128}$, se entrega una cadena no vacía
de 384 bits. Así se ejecuta el \emph{Update} previo.
No se activa el camino de \emph{prediction resistance} que
consume ese argumento en el reseed y lo pone a nulo antes de generar.
Los reseeds necesarios se realizan por separado y se verifica que
se ejecutó la llamada de generación prevista, con $D$ presente
\cite[§9.3.3 y §10.2.1.5.1]{nist}.

\begin{theorem}[Biyectividad a través de $V$]
En la llamada anterior, fijado el estado antiguo $s$,
el mapa entre $A$ y el bloque entregado por CTR\_DRBG es una
permutación de $\bits^{128}$. Preserva exactamente la minentropía
condicional y la distancia a uniforme bajo la misma vista que
en el teorema 1.
\end{theorem}
\begin{proof}
El \emph{Update} previo calcula
$K'=G_K(s)$ y $V'=G_V(s)\oplus A$. La primera
iteración de \emph{Generate} calcula
\[
 Y=\AES_{G_K(s)}\big(\inc(G_V(s)\oplus A)\big).
\]
XOR con una constante, incremento modular y cifrado con clave
fija son permutaciones. Su inversa explícita es
\[
 A=G_V(s)\oplus
       \inc^{-1}\big(\AES^{-1}_{G_K(s)}(Y)\big).
\]
Una vez calculado $Y$, el \emph{Update} posterior usa
$\operatorname{Update}(D,(K',\inc(V')))$. Es posterior a la
computación del bloque y define el estado retenido para la
siguiente llamada, sin cambiar los bits de $Y$. Por eso el
argumento conserva el algoritmo de SP 800-90A Rev. 1,
incluida su actualización posterior. La conservación
estadística sigue de la permutación.
\end{proof}

\paragraph{9-tupla del perfil CTR de un bloque.}
Para expresar este perfil en el mismo formato AIS, se toma
$\mathcal S=\mathcal S_{\rm req}=\bits^{256}\times\bits^{128}$,
$\mathcal R=\bits^{128}$, $\mathcal A=\{\varnothing\}\cup\bits^{128}$,
$I=\{128\}$ y $\phi_0=\operatorname{id}$. Defínanse
$D(\varnothing)=0^{384}$ y $D(a)=0^{256}\|a$ para $|a|=128$.
Las tres funciones restantes son
\begin{align*}
\phi_{\rm req}(s,a)&=
\begin{cases}
s,&a=\varnothing,\\
\operatorname{Update}(D(a),s),&|a|=128,
\end{cases}\\
\psi(k,v)&=\AES_k(\inc(v)),\\
\phi(s,a)&=\operatorname{Update}(D(a),(k,\inc(v))),
\qquad (k,v)=\phi_{\rm req}(s,a).
\end{align*}
La tupla describe el núcleo durante una generación admitida.
La interfaz conserva además el contador de reseed de 90A, rechaza
generación al superar su límite, y sólo libera salida PTG.3 en
la rama $|a|=128$. A efectos de PTG.3.4 el prefijo constante
aporta cero bits frescos. La relación física sigue siendo 128 a 128.

Al recuperar $A$ a partir de $s,Y$, también se conoce $D$ y se puede
reconstruir el siguiente estado con la misma regla. Por tanto, al
fijar la semilla de época y los demás datos auxiliares, la misma
inducción que en la nota principal da una biyección de la secuencia
$A^q\mapsto Y^q$ para peticiones de un bloque. La cota conjunta
conserva la dependencia presente en $A^q$. Los observables de
estado posteriores se incluyen en el modelo de información lateral
si son pertinentes. Este perfil de aplicación usa una llamada
\emph{Generate} estándar sin modificar el algoritmo. Su uso como
postprocesado PTG.3 requiere evaluar PTG.3.2 y la instancia
PTG.2, además de la contabilización de entropía y las pruebas
de salud de la nota principal. La prueba no asigna al perfil una
validación NIST o un claim bajo otro documento de la serie 90.

\subsection{Cuándo se pierde la garantía}

\begin{enumerate}
\item \textbf{Inyectar $A$ de 128 bits directamente.}
El padding NIST sin derivación añade ceros a la derecha. En este caso
$D=A\|0^{256}$, el XOR modifica la mitad superior de la clave
y $V'$ queda fijo. El mapa de interés cruza claves AES para un
texto fijo. La propiedad de permutación de AES no aplica a ese
mapa. No se deduce una colisión concreta de AES a partir de
esa observación, pero la prueba por biyección desaparece.
\item \textbf{Activar la función de derivación.}
\emph{Block\_Cipher\_df} procesa y mezcla la entrada antes del
XOR. La inyección al campo $V$ deja de ser la descrita arriba.
Además, cuando la entrada libre excede 384 bits, una DF de
384 bits no es inyectiva en todo su dominio. La DF requeriría
un modelo independiente para sostener una garantía ITS sobre
la familia restringida de entradas.
\item \textbf{Generar varios bloques con sólo un $A$ fresco.}
Para estado viejo fijo, desde los 128 bits frescos sólo pueden
alcanzarse $2^{128}$ secuencias de salida. Para $q>1$ bloques,
la distancia conjunta a $U_{128q}$ es al menos
$1-2^{-128(q-1)}$. Cada bloque aislado, para una posición fija,
puede seguir siendo una permutación de $A$.
En un llamado desde un estado no condicionado y con un contador
sin vuelta, los bloques de la misma petición son distintos por
la permutación AES. Bajo IID uniforme, el evento de al menos
una repetición tiene probabilidad
$1-\prod_{i=0}^{q-1}(1-i/2^{128})$.
Esta segunda prueba distingue incluso si el estado inicial
fue plenamente uniforme.
\item \textbf{Usar reseed uniforme para un único bloque.}
Existe también una ruta exacta, aunque cuesta 384 bits frescos.
Tras \emph{Reseed} sin DF, si la entrada de entropía de
384 bits es uniforme condicionada al estado anterior y a la
información lateral, $(K',V')$ es conjuntamente uniforme.
Tras ese reseed, una llamada \emph{Generate} de 128 bits
sin additional input emite un bloque uniforme, incluso
condicionado a la nueva clave $K'$. El resultado se extiende
a distancia estadística acotada por procesamiento de datos.
Una cota de minentropía total sobre los 384 bits, por sí
sola, no da independencia entre los campos $K'$ y $V'$ ni
esa exactitud.
\end{enumerate}

\subsection{Conformidad DRG.3 de CTR\_DRBG}

El §102 distingue los 384 bits de estado almacenado y la
dimensión eficaz, que para AES-256 es como máximo 256 bits.
Los perfiles AES-128 y AES-192 quedan fuera de DRG.3 por
dimensión. El §1007 deja explícitamente al solicitante
la prueba para AES-256. La biyección del teorema 3 usa la
variable fresca y el estado viejo fijo. Por tanto, no
resuelve DRG.3.3 ni las propiedades de la evolución
computacional con $D$ conocido o controlado.

Como ejemplo de la obligación computacional, tras un
\emph{Generate} de un bloque, si se conocen el estado
posterior y $D$, eliminar el XOR de la actualización
posterior expone tres evaluaciones consecutivas de AES
bajo la clave de generación. Para enhanced backward secrecy,
hay que justificar por qué esas tres evaluaciones no
permiten predecir eficientemente la evaluación inmediatamente
anterior en el modelo del TOE. Junto con los restantes
requisitos, hace falta analizar la interacción de los estados
y la resistencia a entradas elegidas, incluidas variantes
que sean accesibles fuera del perfil ITS
\cite[§§102, 207, 211--214, 380 y 1007]{ais}\cite[§10.2.1.5.1]{nist}.

\section{Comparación con Hash\_DRBG y HMAC\_DRBG}

En la generación estándar de Hash\_DRBG, el additional input
pasa por un hash junto con el estado y su resultado se combina
con $V$. Se aplica después \emph{Hashgen}. En HMAC\_DRBG,
el additional input entra a una actualización que usa HMAC y
modifica la clave y el valor interno antes de generar.
Las dos composiciones admiten pérdida de información en
dominios amplios y no exponen una permutación estructural
de los 128 bits frescos hacia los 128 bits entregados.
Truncar a 128 bits tampoco da esa propiedad. Se podrían
justificar cotas para distribuciones o construcciones
concretas. De hecho, AIS analiza casos de entropía para
PTG.3 usando modelos de aplicaciones aleatorias.
Eso da una vía de análisis distinta de la igualdad exacta
que se demuestra aquí. En modo de expansión desde una
semilla fija, la limitación de soporte se aplica a los tres
DRBG estándar y al ejemplo de AIS
\cite[§§392--394, 404--413, 1006--1008]{ais}\cite[§§10.1.1--10.1.2]{nist}.

Para una salida de 128 bits con entrada fresca, el núcleo
del §996 usa un cifrado AES y un SHA-256 sobre 384 bits antes
del padding, que cabe en un bloque de compresión SHA-256.
CTR\_DRBG sin DF usa siete cifrados AES por llamada de este
perfil: tres en el \emph{Update} previo, uno de salida y
tres en el posterior. La cuenta excluye siembra,
expansión de claves y control de interfaz. El ejemplo AIS
tiene menos pasos y un estado persistente de 256 bits.
CTR\_DRBG conserva la ventaja práctica de un mecanismo
NIST existente y puede ser atractivo en plataformas
con AES acelerado. La latencia y superficie total de
implementación dependen de la plataforma. Para el dossier
del PTG.3, el núcleo AIS aporta la referencia directa de
§§390(iii) y 996--998. CTR\_DRBG aporta la biyección
exacta de este addendum y deja pendiente la justificación
algorítmica del §1007.

\begin{thebibliography}{9}
\bibitem{ais} M. Peter y W. Schindler.
\emph{A Proposal for Functionality Classes for Random Number Generators}.
BSI, v3.0, 10 de septiembre de 2024.
\href{https://www.bsi.bund.de/SharedDocs/Downloads/EN/BSI/Certification/Interpretations/AIS_31_Functionality_classes_for_random_number_generators_e_2024.pdf?__blob=publicationFile\&v=3}{PDF oficial}.
\bibitem{nist} E. Barker y J. Kelsey.
\emph{Recommendation for Random Number Generation Using Deterministic Random Bit Generators}.
NIST SP 800-90A Rev. 1, junio de 2015.
\url{https://doi.org/10.6028/NIST.SP.800-90Ar1}.
\end{thebibliography}
\end{document}
